\documentclass[11pt]{article}
\usepackage[margin=1.1in]{geometry}
\usepackage{amsmath,amssymb,amsthm,booktabs}
\newtheorem{theorem}{Theorem}[section]
\newtheorem{lemma}[theorem]{Lemma}
\newtheorem{corollary}[theorem]{Corollary}
\newtheorem{conjecture}[theorem]{Conjecture}
\newtheorem{remark}[theorem]{Remark}
\newcommand{\R}{\mathbb{R}}
\newcommand{\Z}{\mathbb{Z}}
\newcommand{\Q}{\mathbb{Q}}
\newcommand{\NC}{N_{C}}
\title{Geometric mass balance on the overlap graph and a two-anchor
reduction of the coincidence conjecture\thanks{Companion section to the
consolidated PSC manuscript (v13 lineage). All spectral claims herein are
conditional content in the sense of v13: nothing in this note proves the Pisot
substitution conjecture; it reorganizes the open core.}}
\author{}
\date{June 12, 2026 (v9 cross-reference revision, September 8, 2026)}
\begin{document}
\maketitle

\section{Setup}

Let $\sigma$ be a primitive irreducible Pisot substitution on
$\mathcal{A}=\{1,\dots,k\}$ with incidence data $A_{ab}=|\sigma(a)|_b$,
Perron value $\beta>1$ of degree $k$, and tile lengths
$L=(L_1,\dots,L_k)>0$ with $AL=\beta L$. Irreducibility of the
characteristic polynomial makes $L_1,\dots,L_k$ linearly independent over
$\Q$. A \emph{genuine overlap} is a triple $s=(i,\delta,j)$ with
$\delta\in\Z^k$, $t(s)=\langle\delta,L\rangle$, such that the tiles
$[0,L_i)$ and $[t,t+L_j)$ intersect with nonempty interior; write
$o(s)$ for the intersection and $w(s)=|o(s)|>0$. Inflation sends $s$ to
the cells of the common refinement of the $\sigma(i)$- and
$\sigma(j)$-subdivisions restricted to $\beta\cdot o(s)$; each cell is a
genuine overlap $(i',\delta',j')$ with
$\delta'=A^{\!\top}\delta+\mathrm{pp}_i(p)-\mathrm{pp}_j(q)$ for
prefix-Parikh vectors $\mathrm{pp}$. A \emph{coincidence} is a state
$(i,0,i)$. Overlap coincidence is equivalent to pure discreteness of the
suspension tiling dynamics \cite{Sol97,AL11}. The state space reachable
from return-vector seeds is finite (Meyer-set arguments; cf.\
\cite{AL11}), so no balanced-pair termination input is used anywhere
below.

\section{Mass balance}

\begin{theorem}[Partition identity]\label{thm:partition}
For every genuine overlap $s$, counting children with multiplicity,
$\beta\,w(s)=\sum_{c\,\mathrm{child\,of}\,s}w(c)$.
\end{theorem}

\begin{proof}
$\beta\cdot[0,L_i)$ is partitioned (up to endpoints) by the sub-tile
intervals of $\sigma(i)$, and $\beta\cdot[t,t+L_j)$ by those of
$\sigma(j)$. The common refinement of the two partitions, restricted to
$\beta\cdot o(s)$, is an interval partition of $\beta\cdot o(s)$ whose
cells of nonempty interior are exactly the child overlaps. Sum lengths.
\end{proof}

Let $C$ be a recurrent noncoincident strongly connected component (SCC)
of the genuine overlap graph, $\NC$ its internal child-count matrix
(nonnegative, integer, irreducible), $w|_C>0$, and
$\mathrm{leak}(s)\ge 0$ the total measure of children of $s$ outside
$C$. Theorem~\ref{thm:partition} restricts to
\begin{equation}\label{eq:mb}
\beta\,w \;=\; \NC\,w+\mathrm{leak}\qquad\text{on }C.
\end{equation}

\begin{theorem}[Ceiling and leak dichotomy]\label{thm:dichotomy}
$\rho(\NC)\le\beta$, with equality iff $\mathrm{leak}\equiv 0$ on $C$.
\end{theorem}

\begin{proof}
$\NC w\le\beta w$ entrywise with $w>0$ gives $\rho(\NC)\le\beta$
(Collatz--Wielandt). If $\mathrm{leak}\equiv0$ then $w$ is a positive
eigenvector, hence the Perron eigenvector, and $\rho(\NC)=\beta$. If
$\mathrm{leak}\not\equiv0$, pair \eqref{eq:mb} with the left Perron
eigenvector $u>0$ of $\NC$:
$(\beta-\rho(\NC))\langle u,w\rangle=\langle u,\mathrm{leak}\rangle>0$.
\end{proof}

\begin{theorem}[Exact gap formula]\label{thm:gap}
$\displaystyle \beta-\rho(\NC)
 \;=\;\frac{\langle u,\mathrm{leak}\rangle}{\langle u,w\rangle}.$
\end{theorem}

\begin{remark}
Theorem~\ref{thm:partition} is elementary and implicit in the framework
of \cite{AL11}; we record it because \eqref{eq:mb} and
Theorems~\ref{thm:dichotomy}--\ref{thm:gap} follow from it in two lines
each, and because the identical pairing applied to the balanced-pair
mass-balance identity of the consolidated manuscript yields the same gap
formula on the symbolic side. The spectral gap of a recurrent overlap
SCC \emph{is} its eigenvector-weighted coincidence-production rate. For
a \emph{closed} $C$ (no noncoincident child outside $C$), leak is
coincidence mass only, so $\rho(\NC)<\beta$ iff some state of $C$ has a
coincidence child. The open core of the conjecture is therefore the same
statement on both the symbolic and geometric sides: \emph{a closed
recurrent noncoincident component must leak}.
\end{remark}

\begin{theorem}[Size floor for the counterexample shape]\label{thm:floor}
If $C$ is closed, recurrent, noncoincident with $\mathrm{leak}\equiv0$,
then $\beta\in\mathrm{spec}(\NC)$, hence the characteristic polynomial
of $A$ divides that of $\NC$ and $|C|\ge k$. Moreover, writing
$w=VL$ with the integer matrix $V$ determined by the four overlap cases
$v_s\in\{e_j,\,e_i-\delta_s,\,\delta_s+e_j,\,e_i\}$, linear independence
of the $L_a$ upgrades \eqref{eq:mb} to the integer intertwining
$\NC V=VA$ with $V$ injective.
\end{theorem}

\begin{proof}
$\NC w=\beta w$ with $w\ne0$, $w\in\Q(\beta)^{|C|}$, exhibits $\beta$ as
an eigenvalue of an integer matrix; the minimal polynomial of $\beta$ is
the (irreducible) characteristic polynomial of $A$, of degree $k$.
Independence of the $L_a$ converts $\NC VL=VAL$ into $\NC V=VA$; since
$\Q^k$ is a simple $\Q[A]$-module, $V\ne0$ is injective.
\end{proof}

\section{The two anchors}

Call $s=(i,\delta,j)$ \emph{prefix-aligned} if $\delta=0$ (the tiles
share their left endpoint) and \emph{suffix-aligned} if
$\delta=e_i-e_j$ (then $t=L_i-L_j$: the tiles share their right
endpoint). The prefix-aligned children of a prefix-aligned state are the
equal-prefix-Parikh cells (the Barge--Diamond graph
\cite{BD02,AI01}); the suffix-aligned children of a suffix-aligned
state are the equal-suffix-Parikh cells, i.e.\ the Barge--Diamond graph
of the reversed substitution $\tilde\sigma$ (same incidence matrix).

\begin{theorem}[Prefix channel]\label{thm:B}
Let $C$ be closed, recurrent, noncoincident, containing a prefix-aligned
state $(i,0,j)$. If $(i,j)$ satisfies strong coincidence (some
$\sigma^n(i)$, $\sigma^n(j)$ share a letter at equal prefix-Parikh),
then $C$ produces a coincidence child.
\end{theorem}

\begin{theorem}[Suffix channel]\label{thm:Bp}
Same statement with a suffix-aligned anchor $(i,e_i-e_j,j)$ and strong
coincidence at equal \emph{suffix}-Parikh (equivalently, strong
coincidence of $\tilde\sigma$ for $(i,j)$).
\end{theorem}

\begin{proof}[Proof of both]
Anchor the two inflation towers at the common endpoint. For each level
$m\le n$, the sub-tile pairs containing the depth-$n$ coincidence tile
both contain its interior, hence form genuine overlaps; by closedness
this nested chain lies in $C$ until its final step exhibits a
coincidence child of a state of $C$.
\end{proof}

\begin{remark}[$\gamma$-channel]\label{rem:gchannel}
The proofs of Theorems~\ref{thm:B}--\ref{thm:Bp} use only that the two
towers are anchored at a common point realized by matching control
points. The same nested-chain argument proves: if a closed recurrent
noncoincident SCC contains a state $(i,\delta,j)$ with
$\langle\delta,L\rangle=c_j-c_i$ for an admissible control-point system
$(c_a)$ arising from a tile map of some power of the substitution, and
$(i,j)$ satisfies strong coincidence with respect to that system, then
the SCC produces. Theorems~\ref{thm:B} and~\ref{thm:Bp} are the
instances given by the first-letter and last-letter tile maps; we
single them out because Conjecture~\ref{conj:H3} asserts that these two
anchors always occur, whereas \cite{Ak16} manufactures the anchor per
trapped component (Remark~\ref{rem:calib}).
\end{remark}

\begin{lemma}[Cycle pinning]\label{lem:pin}
A directed cycle of the overlap graph with edge digit word
$(d_1,\dots,d_n)\subset\Z^k$ carries the uniquely determined integer
offset orbit
$\delta_1=(I-(A^{\!\top})^n)^{-1}\sum_{i=1}^{n}(A^{\!\top})^{n-i}d_i$,
since $\det(I-(A^{\!\top})^n)=\prod_m(1-\beta_m^n)\neq0$ for Pisot
$\beta$.
\end{lemma}

\begin{lemma}[Shared cut]\label{lem:cut}
Let $s$ belong to a closed recurrent noncoincident SCC $C$ and suppose a
refinement window of $s$ contains an interior cut point common to both
grids. Then $C$ contains a prefix- or suffix-aligned state.
\end{lemma}

\begin{proof}
At a common cut both grids end one sub-tile and begin another; the cell
ending there is right-aligned, the cell beginning there is left-aligned.
A coincident cell is both, and its endpoints are again common cuts, so
the maximal coincidence run through the cut is bounded by common cuts
whose outside-adjacent cells are aligned and noncoincident, hence in $C$
by closedness. If the run is flush against one window edge, the aligned
noncoincident cell at its boundary on the \emph{other} side serves; a
run flush against both edges fills the window and leaves $s$ without
noncoincident children, contradicting recurrence.
\end{proof}

\begin{theorem}[Anchors from production]\label{thm:C}
Every producing closed recurrent noncoincident SCC contains a prefix- or
suffix-aligned state.
\end{theorem}

\begin{proof}
A coincidence cell is aligned on both sides, so its endpoints are common
cuts; apply Lemma~\ref{lem:cut}.
\end{proof}

By Theorem~\ref{thm:C}, anchoring is automatic for producing components;
the conjectural content concerns only the trapped shape of
Theorems~\ref{thm:dichotomy} and~\ref{thm:floor}:

\begin{conjecture}[Trapped anchoring, H3$'$]\label{conj:H3}
Every closed recurrent noncoincident SCC with $\mathrm{leak}\equiv0$
contains a prefix-aligned or a suffix-aligned state.
\end{conjecture}

In a $\mathrm{leak}\equiv0$ component no coincidence cells exist, so by
Lemma~\ref{lem:cut} \emph{any} common interior cut anchors it
immediately. Under the hypotheses of
Theorem~\ref{thm:reduction} a surviving counterexample must therefore
keep its two cut grids totally disjoint at every refinement level, in
windows growing like $\beta^n$ --- an alignment-avoidance condition at
all scales, layered on the size floor and intertwining of
Theorem~\ref{thm:floor}. The Pisot property places $\beta^n t$ at
bounded distance from the cut module; ``bounded'' versus ``zero'' is the
residue, in contact with Property~(F)-type finiteness statements of
$\beta$-numeration.

\begin{theorem}[Two-anchor reduction]\label{thm:reduction}
If $\sigma$ and its reversal $\tilde\sigma$ satisfy the strong
coincidence condition and Conjecture~\ref{conj:H3} (trapped anchoring) holds for $\sigma$,
then every closed recurrent noncoincident SCC produces; consequently
overlap coincidence holds and the tiling dynamical system of $\sigma$
has pure discrete spectrum.
\end{theorem}

\begin{proof}
Theorems~\ref{thm:B}--\ref{thm:Bp} on the anchor supplied by
Conjecture~\ref{conj:H3}; every closed recurrent SCC leaks, and
nonclosed recurrent SCCs leak by definition through their escapes
together with Theorem~\ref{thm:dichotomy}; conclude by
\cite{Sol97,AL11}.
\end{proof}

\begin{remark}[Calibration against the literature]\label{rem:calib}
Akiyama--Lee \cite{AL14} prove overlap coincidence $\Rightarrow$
(generalized) strong coincidence when the height group is trivial,
which holds for one-dimensional irreducible Pisot substitutions
\cite{BK06,Sing06}. Akiyama \cite[Thm.~3.2, Cor.~3.3]{Ak16} proves:
\emph{multiple strong coincidence of level $n$} (strong coincidence for
all admissible control-point systems of $\Omega^n$ with
$\Lambda-\Lambda\subset\langle G\rangle$) is equivalent to overlap
coincidence for the suspension tiling, with $n$ bounded by the number
of overlap classes. His proof shows the hypothetical trapped component
has spectral radius $|\det Q|=\beta$ --- precisely the
$\mathrm{leak}\equiv0$ shape of Theorems~\ref{thm:dichotomy}
and~\ref{thm:floor} --- and then \emph{constructs} a tile map
(\cite[Lem.~3.5]{Ak16}) whose control points match the offset of a
mixed-letter state of the trapped component: the anchor is moved to
the offset. Thus $\gamma$-generalized anchoring of the trapped
component is unconditional in \cite{Ak16}, and the residue of his
reduction is strong coincidence for the manufactured systems.
Prefix- and suffix-strong coincidence are the control-point systems of
the first-letter and last-letter tile maps (by
Lemma~\ref{lem:pin}, these are exactly the tile maps whose cycles carry
zero digit corrections, pinning $\delta$ to the two endpoint anchors).
Theorem~\ref{thm:reduction} is therefore a sharpening proposal for the
minimization question of \cite[Rem.~3.4]{Ak16} (``can we take
$n=1$?''): conditional on Conjecture~\ref{conj:H3}, the two endpoint
systems suffice. Conjecture~\ref{conj:H3} is strictly stronger than
what \cite{Ak16} needs --- its failure shape, a closed component whose
pinned offsets avoid $\{0\}\cup\{L_i-L_j\}$, is exactly what his
construction routes around, and exactly what the censuses below never
exhibit.
\end{remark}

\section{The $\ell$-fixed reduction of Conjecture~\ref{conj:H3}}

For a state of a $\mathrm{leak}\equiv0$ component every refinement cell
is noncoincident and in the component, so the leftmost- and
rightmost-cell maps are self-maps of the component; being eventually
periodic on a finite set, \emph{every trapped component contains an
$\ell$-cycle and an $r$-cycle}, and a period-$p$ $\ell$-cycle state is
$\ell$-fixed for $\sigma^p$ (the leftmost cell of the level-$p$
refinement is the $p$-fold leftmost cell). Trappedness persists under
powers and descendant cones are power-independent.

\begin{lemma}[Telescoping]\label{lem:tel}
Let $s^*=(i^*,\delta^*,j^*)$ be $\ell$-fixed with $t^*>0$, so that
$\sigma(j^*)$ begins with $j^*$, $\sigma(i^*)=u^*\,i^*\,v^*$, and
$(A^{\!\top}-I)\delta^*=\pi(u^*)$ (the matrix is invertible since $1$ is
not an eigenvalue of $A$; $v^*\neq\varepsilon$ is forced by
$t^*<L_{i^*}$). Then, exactly,
$(A^{\!\top})^n\delta^*=\delta^*+\pi(Q_n)$ where
$Q_n=\sigma^{n-1}(u^*)\cdots\sigma(u^*)u^*$ is the prefix of
$\sigma^n(i^*)$ preceding the tracked central occurrence of $i^*$.
\end{lemma}

Consequently a level-$n$ shared cut for $s^*$ is the \emph{static}
membership $\delta^*\in\{\pi(a)-\pi(v)\}$ over factors $a$ anchored at
the central occurrence and prefixes $v$ of $\sigma^n(j^*)$; in the limit,
over prefixes of the two $\sigma$-fixed half-tilings
$\omega_I=i^*v^*\sigma(v^*)\cdots$ and $\omega_J=\lim\sigma^n(j^*)$ at
offset $t^*$. This is the algebraic coincidence of Lee \cite{Lee07}
localized at the single vector $\delta^*$; an unconditional proof would
prove the Pisot conjecture.

\begin{theorem}[$\ell$-fixed reduction]\label{thm:lfix}
Conjecture~\ref{conj:H3} holds if for every $p\ge1$ the finite statement
$Q(p)$ holds: every in-window, integral, $\ell$- or $r$-fixed state of
$\sigma^p$ contains an anchor or a coincidence in its descendant cone.
\end{theorem}

Two unconditional fences accompany this: $\delta^*=(A^{\!\top}-I)^{-1}
\pi(u^*)$ must be \emph{integral} to host a state at all (the fence
removes $20\%$, $55\%$, $85\%$, $97\%$ of candidate data at
$p=1,2,3,6$ in the censuses below), and nearest-vertex difference
vectors have bounded Perron coordinate and bounded conjugates, hence
form a finite set per specimen.

\begin{lemma}[Norm floor]\label{lem:floor}
An integral $\ell$-fixed state of any power $\sigma^p$ with
$\delta^*\neq0$ satisfies
$|\langle\delta^*,L^{(m)}\rangle|\le R'_m:=B_m/(1-|\beta_m|)^2$
uniformly in $p$ (Dumont--Thomas prefix bound \cite{DT89} over
$|\beta_m^p-1|\ge1-|\beta_m|$), and
$|t^*|\ge\kappa(\sigma):=1/(D^k\prod_{m\ge2}R'_m)>0$, where $D$ clears
the denominators of $L$, since the field norm of $t^*$ is a nonzero
rational of bounded denominator. In particular no genuine integral
state has offset in $(0,\kappa)$: the small-offset regime is empty, and
the integrality fence observed in the censuses is this floor in action.
\end{lemma}

\begin{theorem}[Finite equivalence]\label{thm:qinf}
Let $\Delta(\sigma)=\{\delta\in\Z^k\setminus\{0\}:
|\langle\delta,L\rangle|<L_{\max},\
|\langle\delta,L^{(m)}\rangle|\le R'_m\}$, a finite explicit set
containing every $\ell$- or $r$-fixed offset of every power. If
$Q_\infty(\sigma)$ holds --- \emph{every genuine state
$(i,\delta,j)$ with $\delta\in\Delta(\sigma)$ contains an anchor or a
coincidence in its (finite) descendant cone} --- then
Conjecture~\ref{conj:H3} holds for $\sigma$; conversely, failure of
$Q_\infty$ exhibits an anchor-free trapped component, refuting
Conjecture~\ref{conj:H3}. Thus $Q_\infty(\sigma)$ is \emph{equivalent}
to trapped anchoring, and the quantifier over powers in
Theorem~\ref{thm:lfix} is eliminated: Conjecture~\ref{conj:H3} is a
single finite, effective check per substitution.
\end{theorem}

\begin{proof}
($\Rightarrow$) A trapped component contains an $\ell$-cycle,
$\ell$-fixed for $\sigma^p$; its offset is integral, conjugate-bounded
and norm-floored by Lemma~\ref{lem:floor}, hence in $\Delta(\sigma)$.
$Q_\infty$ places an anchor or coincidence in its cone, which lies in
$C\cup\{\text{coincidences}\}$ by closedness; a coincidence contradicts
$\mathrm{leak}\equiv0$. ($\Leftarrow$) If some $\Delta$-state's cone
avoids anchors and coincidences, the cone is finite (conjugates contract
under inflation with bounded corrections), every state has a child, and
a bottom SCC of its condensation has all children inside itself: closed,
recurrent, with $\mathrm{leak}\equiv0$ and no anchor --- an anchor-free
trapped component. Note also that such a component is \emph{totally
cut-disjoint unconditionally}: by Lemma~\ref{lem:cut} with empty
coincidence runs, any common interior cut would place a noncoincident
aligned cell in it.
\end{proof}

Per-specimen certificates (tail-corrected radius; an earlier
under-converged radius computation was caught in audit and its outputs
superseded, and a subsequent exact tail-corrected suffix radius---the
superseding instrument---shrinks $\Delta$ a further $\sim2.5\times$ with
identical verdicts, e.g.\ flipped Tribonacci $|\Delta|=402\to144$,
plastic $2{,}504\to1{,}134$, re-certified states $2{,}868\to1{,}050$ and
$17{,}460\to7{,}920$ respectively): $Q_\infty$ is verified for the
\emph{complete exhaustive} $k=3$, $|\sigma(a)|\le3$ corpus --- all
$4{,}554$ primitive irreducible Pisot specimens of the $59{,}319$-index
enumeration, coverage confirmed by set equality against the directly
enumerated PIP index family (a $47$-index shard hole inside previously
claimed coverage was exposed by exactly this reconciliation and
gap-filled) --- with $0$ failures, $0$ capped, $0$ timeouts. The deepest
specimens were closed by the optimal-slab and core certificates of
\S\ref{sec:core} (largest retained closure $193{,}563$ states; the
former multi-million-state band collapses under the core slab). A
further $24$ random $k=4$ specimens
($|\sigma(a)|\le3$, $|\beta_2|$ up to $0.976$) certify with $0$
failures, the five deepest via the core certificate of
Theorem~\ref{thm:corered} (closures to $11.4$ million states).
Conjecture~\ref{conj:H3} therefore holds unconditionally for each
certified specimen. Per-specimen decidability of pure
discreteness is of course already given by \cite{AL11}; the content
here is the uniform target: the conjecture now reads ``$Q_\infty(\sigma)$
for all primitive irreducible Pisot $\sigma$'', a statement about a
norm-floor-bounded finite offset family, with no powers and no
dynamics over $p$.

\section{The saturated core}\label{sec:core}

For $m\ge2$ let $B_m$ denote the one-step prefix-conjugate bound of
Lemma~\ref{lem:floor},
\[
B_m \;=\; \max_{a\in\mathcal{A}}\ \max_{p\ \text{prefix of}\
\sigma(a)}\ \bigl|\langle\pi(p),L^{(m)}\rangle\bigr|,
\]
and for $\varepsilon\ge0$ define the \emph{saturated core}
\[
\mathcal{C}_\varepsilon(\sigma)\;=\;\Bigl\{\,\text{genuine }(i,\delta,j)
\;:\;\bigl|\langle\delta,L^{(m)}\rangle\bigr|\;\le\;
\frac{2B_m}{1-|\beta_m|}+\varepsilon\ \ \text{for all }m\ge2\,\Bigr\}.
\]

\begin{lemma}[Core invariance]\label{lem:coreinv}
$\mathcal{C}_\varepsilon$ is forward-invariant under overlap inflation
for every $\varepsilon\ge0$.
\end{lemma}

\begin{proof}
A child offset is $\delta'=A^{\!\top}\delta\pm
(\mathrm{pp\ difference})$, so
$\langle\delta',L^{(m)}\rangle=\beta_m\langle\delta,L^{(m)}\rangle+
\rho$ with $|\rho|\le2B_m$. Writing $F_m=2B_m/(1-|\beta_m|)$: if
$|\langle\delta,L^{(m)}\rangle|\le F_m+\varepsilon$ then
$|\langle\delta',L^{(m)}\rangle|\le|\beta_m|(F_m+\varepsilon)+2B_m=
F_m+|\beta_m|\varepsilon\le F_m+\varepsilon$.
\end{proof}

\begin{lemma}[Entry]\label{lem:entry}
Let $\varepsilon>0$ and let $S$ be any family of genuine states with
$|\langle\delta,L^{(m)}\rangle|\le R_m$ for all $m$. Every descendant
chain of every $s\in S$ enters $\mathcal{C}_\varepsilon$ within
$h(\varepsilon)=\max_m\lceil\log((R_m-F_m)/\varepsilon)/
\log(1/|\beta_m|)\rceil$ inflation levels.
\end{lemma}

\begin{proof}
Every genuine state has at least one child by
Theorem~\ref{thm:partition} ($w(s)>0$). Along any chain the recursion of
Lemma~\ref{lem:coreinv} gives
$|\langle\delta_n,L^{(m)}\rangle|-F_m\le|\beta_m|^n(R_m-F_m)$.
\end{proof}

\begin{lemma}[Recurrent states are core]\label{lem:reccore}
Every state lying on a directed cycle of the genuine overlap graph
belongs to $\mathcal{C}_0$ (the \emph{closed} ball; no $\varepsilon$).
\end{lemma}

\begin{proof}
If $s$ lies on a cycle of length $c$, it is its own $nc$-step
descendant for every $n$; iterating the recursion around the cycle,
$|\langle\delta_s,L^{(m)}\rangle|\,(1-|\beta_m|^{nc})\le
2B_m\,(1-|\beta_m|^{nc})/(1-|\beta_m|)$; divide.
\end{proof}

\begin{theorem}[Core reduction]\label{thm:corered}
For every $\varepsilon>0$ the \emph{core certificate}
$Q^{\mathrm{core}}_\infty(\sigma)$ --- every state of
$\mathcal{C}_\varepsilon(\sigma)$ contains an anchor or a coincidence
in its descendant cone --- is equivalent to trapped anchoring
(Conjecture~\ref{conj:H3} for $\sigma$), and the cone may be searched
inside $\mathcal{C}_\varepsilon$ alone: anchors and coincidences lie in
$\mathcal{C}_0$ (for $\delta=0$ the conjugates vanish; for
$\delta=e_i-e_j$ they are bounded by $2B_m$), and
$\mathcal{C}_\varepsilon$ is invariant, so the certificate is
self-contained.
\end{theorem}

\begin{proof}
If a trapped component exists, its states are recurrent, hence in
$\mathcal{C}_0\subseteq\mathcal{C}_\varepsilon$ by
Lemma~\ref{lem:reccore}; the core certificate places an anchor or
coincidence in such a state's cone, which lies in
$C\cup\{\text{coincidences}\}$ by closedness --- either contradicts
$\mathrm{leak}\equiv0$ exactly as in Theorem~\ref{thm:qinf}.
Conversely, if some $\mathcal{C}_\varepsilon$-state's cone avoids
anchors and coincidences, the cone is finite (its conjugates are
bounded by invariance and its Perron coordinates by the window), every
state has a child, and a bottom SCC of its condensation is an
anchor-free trapped component, as in Theorem~\ref{thm:qinf}.
\end{proof}

\begin{remark}[Optimal slab; norm-floor upgrade]\label{rem:optslab}
Lemma~\ref{lem:reccore} pins every cycle state --- in particular the
$\ell$-cycle witness of any trapped component --- to conjugate radius
$2B_m/(1-|\beta_m|)$, improving the bound of Lemma~\ref{lem:floor} from
quadratic to linear in $(1-|\beta_m|)^{-1}$ whenever $|\beta_m|>1/2$,
with the offset floor $\kappa(\sigma)$ rising correspondingly. Neither
bound dominates globally ($B_m/(1-|\beta_m|)^2$ is smaller for
$|\beta_m|<1/2$); the optimal certificate slab takes
$\min\bigl(B_m/(1-|\beta_m|)^2,\;2B_m/(1-|\beta_m|)\bigr)$ per
embedding. Per cycle, the pinned orbit of Lemma~\ref{lem:pin}
satisfies the exact refinement
$|\langle\delta_s,L^{(m)}\rangle|\le
2B_m\sum_{i<c}|\beta_m|^{\,i}/\,|1-\beta_m^{\,c}|$; for complex
$\beta_m$ away from resonance ($\arg\beta_m^{\,c}$ away from $0$) this
is $O(B_m\,c)$, and the Pisot property excludes exact resonance
($\prod_m(1-\beta_m^{\,n})\ne0$, Lemma~\ref{lem:pin}) but not
near-resonance, so no uniform $O(B_m)$ bound follows; the
near-resonant regime is the identified hard case. Conjecture~\ref{conj:H3}
in its sharpest form now reads: \emph{every recurrent noncoincident
state with conjugates at the one-step-saturated scale
$2B_m/(1-|\beta_m|)$ lies in a producing component.}
\end{remark}

\begin{remark}[Computational verification; the core as enabling step]
\label{rem:coreverif}
Exact verification (no tolerance beyond $10^{-9}$ slab fattening):
core invariance holds with zero escapees on all $29$ computed core
closures ($5$ at $k=3$, $24$ at $k=4$); on the largest $k=3$ specimen
the closed-ball bound is realized to four decimals (maximal conjugate
$39.1870$ against radius $39.1880$). Lemma~\ref{lem:reccore} verified
on all $24$ evaluable closures ($k=3$ and $k=4$): all recurrent states
inside the core in every embedding. The recurrent skeleton is
consistently minuscule ($0.01$--$2.9\%$ of the closure; $557$ of
$600{,}357$ states on a representative deep specimen). The reduction
factor of the core certificate scales like $(1-|\beta_2|)^{-2}$ and
reaches $87\times$ at $k=3$, $|\beta_2|=0.94$; at $k=4$ it is an
\emph{enabling} step rather than an optimization --- the
$R'$-slab certificate exhausts a $4$\,GB working set at
$|\beta_2|\approx0.946$ while the core certificate completes at
$|\beta_2|=0.976$ with an $11{,}434{,}672$-state closure. As throughout,
these verifications are elimination, not validation: no census bears on
the uniform statement.
\end{remark}

\section{Powers, period, and the trapped shape}\label{sec:power}

\begin{lemma}[Pisot field identity]\label{lem:field}
For every Pisot number $\beta$ and every $p\ge1$,
$\mathbb{Q}(\beta^p)=\mathbb{Q}(\beta)$. Consequently the $k$ values
$\beta_m^{\,p}$ are pairwise distinct, and the characteristic polynomial
of $A^p$ is the (irreducible) minimal polynomial of $\beta^p$.
\end{lemma}

\begin{proof}
If an embedding $\sigma$ fixes $\beta^p$ but moves $\beta$, then
$\sigma(\beta)=\beta\zeta$ with $\zeta^p=1$, $\zeta\ne1$, so
$|\sigma(\beta)|=\beta>1$, contradicting the Pisot property. Hence
$\deg\beta^p=k$, so the Galois orbit of $\beta^p$ has $k$ distinct
elements; these lie among the $k$ values $\beta_m^{\,p}$, which are
therefore pairwise distinct, and
$\chi_{A^p}=\prod_m(x-\beta_m^{\,p})$ is the minimal polynomial of
$\beta^p$.
\end{proof}

\begin{corollary}[Power closure of the PIP class]\label{cor:power}
If $\sigma$ is primitive irreducible Pisot, so is $\sigma^p$ for every
$p\ge1$. Moreover, when $\beta$ has a complex conjugate pair,
$\arg\beta_2\notin\pi\mathbb{Q}$: a rational angle would give
$\beta_2^{\,2b}=\beta_3^{\,2b}$ for some $b$, collapsing the Galois
orbit of $\beta^{2b}$ below degree $k$.
\end{corollary}

\begin{theorem}[Period-scaled size floor]\label{thm:periodfloor}
Let $C$ be a closed recurrent noncoincident component with
$\mathrm{leak}\equiv0$ and period $p$, with the integer intertwining
$N_CV=VA$, $V$ injective, of Theorem~\ref{thm:floor}. Write
$V_t=P_tV$ for the coordinate projection $P_t$ onto the cyclic class
$C_t$. Then $N_CV_t=V_{t+1}A$ (indices mod $p$), every $V_t$ is
injective, $|C_t|\ge k$ for every $t$, and $|C|\ge p\,k$. The
mass-balance eigen-relation refines classwise:
$N_C\,w_t=\beta\,w_{t+1}$ with $w_t=V_tL>0$.
\end{theorem}

\begin{proof}
$N_C$ maps class $t$ into class $t+1$, so projecting $N_CV=VA$ gives
$N_CV_t=V_{t+1}A$. If some $V_t=0$ then $V_{t+1}A=0$ and $A$ is
invertible, so $V_{t+1}=0$; around the cycle $V=0$, a contradiction.
Composing $p$ steps, $N_C^{\,p}V_t=V_tA^p$, so $\ker V_t$ is
$A^p$-invariant. By Lemma~\ref{lem:field}, $\mathbb{Q}[A^p]$ acts on
$\mathbb{Q}^k\cong\mathbb{Q}(\beta)$ as multiplication by $\beta^p$
with $\mathbb{Q}(\beta^p)=\mathbb{Q}(\beta)$, so $\mathbb{Q}^k$ is a
simple $\mathbb{Q}[A^p]$-module; $V_t\ne0$ forces $V_t$ injective.
\end{proof}

\begin{theorem}[Power reduction; aperiodicity without loss of
generality]\label{thm:powerred}
Let $C$ be a trapped component of period $p$ for $\sigma$. Then each
cyclic class $C_t$ is a trapped component for $\sigma^p$ --- closed
(all $p$-step descendants of $C_t$ remain in $C_t$), recurrent with
\emph{primitive} internal matrix $B_t=N_C^{\,p}|_{C_t}$, noncoincident,
$\mathrm{leak}\equiv0$, and anchor-free if $C$ is. Since $\sigma^p$ is
again primitive irreducible Pisot (Corollary~\ref{cor:power}),
Conjecture~\ref{conj:H3} for all primitive irreducible Pisot
substitutions is equivalent to its restriction to trapped components
with primitive (aperiodic) internal matrix. Every attack on, or hunt
for a counterexample to, uniform trapped anchoring may therefore assume
$N_C$ primitive.
\end{theorem}

\begin{proof}
Closedness with $\mathrm{leak}\equiv0$ means every one-step descendant
of $C$ lies in $C$; $p$-step descendants of $C_t$ lie in
$C_{t+p}=C_t$. $B_t$ is irreducible and aperiodic by the
Perron--Frobenius cyclic-class theory; the remaining properties are
inherited state-by-state, and the anchor set
$\{\delta=0\}\cup\{\delta=e_i-e_j\}$ does not depend on the
substitution. Conversely, an aperiodic trapped component is a trapped
component, so the restricted conjecture implies nothing weaker.
\end{proof}

\begin{theorem}[Frozen rays]\label{thm:frozen}
Up to replacing $\sigma$ by a power (preserving the primitive irreducible
Pisot class by Corollary~\ref{cor:power} and trappedness with primitive
$N_C$ by Theorem~\ref{thm:powerred}), every trapped component contains
an $\ell$-fixed state $s^{*}_{\ell}$ and an $r$-fixed state
$s^{*}_{r}$, both with $t^{*}\ne0$. For each, the associated pair of
$\sigma$-fixed half-tilings at offset $t^{*}$ (the setting of
Lemma~\ref{lem:tel}) is cut-disjoint along the \emph{entire} ray:
no level ever produces a common cut.
\end{theorem}

\begin{proof}
The leftmost-cell map $\ell$ is a self-map of $C$ (closedness with
$\mathrm{leak}\equiv0$), hence has a cycle of some length $q_\ell$; a
cycle state is $\ell$-fixed for $\sigma^{q_\ell}$, and likewise for
$r$. Pass to $\sigma^{q}$, $q=\mathrm{lcm}(q_\ell,q_r)$: this is again
primitive irreducible Pisot, $C$ is again trapped with
$N_C^{\,q}$ primitive, and the chosen cycle states are $\ell$- and
$r$-fixed. Anchor-freeness excludes $\delta=0$, and
$\langle\delta,L\rangle=0$ iff $\delta=0$ by linear independence of the
$L_a$, so $t^{*}\ne0$. A common cut at any refinement level places, by
Lemma~\ref{lem:cut}, an aligned (anchor) state in $C$, contradicting
anchor-freeness.
\end{proof}

\begin{remark}[Honest placement]\label{rem:frozenplacement}
Theorem~\ref{thm:frozen} localizes the open core to two concrete
infinite objects per hypothetical counterexample: two pairs of
$\sigma$-fixed rays, offset by core-bounded $t^{*}_{\ell},t^{*}_{r}$,
whose cut sets are disjoint forever --- equivalently, whose
prefix-Parikh difference walk, confined to the finite saturated core of
\S\ref{sec:core}, avoids the fixed lattice point $\delta^{*}$ at every
step. This is Lee's algebraic coincidence localized as in
Lemma~\ref{lem:tel}, now with the avoidance walk's range pinned
to the core and the substitution normalized to be aperiodic with
$\ell$- and $r$-fixed witnesses. It is a sharpening of the
counterexample shape, not progress on the avoidance question itself:
forever avoiding one point of a finite visited set is a positional
constraint that no count-level or marginal argument can refute.
\end{remark}

\begin{theorem}[Transversal substitution]\label{thm:transversal}
Let $C$ be a trapped component, normalized as in
Theorem~\ref{thm:frozen} ($N_C$ primitive, $s^{*}_{\ell}$
$\ell$-fixed). The map $\Sigma_C\colon s\mapsto(\text{children of }s
\text{ in left-to-right order})$ is a well-defined substitution on the
finite alphabet $C$ (trappedness excludes coincidence letters), with
abelianization $N_C$, primitive, admitting the one-sided fixed point
$\tau=\lim\Sigma_C^{\,n}(s^{*}_{\ell})$; and $\tau$ is exactly the
transversal sequence of the frozen ray pair of
Theorem~\ref{thm:frozen} --- the left-to-right sequence of overlap
cells of the two rays' common refinement, every cell of which lies in
$C$. The natural tile lengths $w(s)$ realize $\tau$ as a self-similar
tiling of the ray with inflation $\beta$. Consequently a trapped
component is equivalent to the following package: a primitive
substitution $\Sigma$ on $|C|\ge k$ letters with reducible
characteristic polynomial $\chi_A\cdot q(x)$ and Pisot--Perron
expansion $\beta$ of algebraic degree $k<|C|$, equipped with the
integer valuation $s\mapsto v_s$ intertwining $N_C$ with $A$
(Theorem~\ref{thm:floor}), whose fixed point never meets the
coincidence diagonal.
\end{theorem}

\begin{proof}
Children of a state, being the cells of an interval partition of
$\beta\cdot o(s)$, carry a canonical left-to-right order, and by
closedness with $\mathrm{leak}\equiv0$ all lie in $C$, so $\Sigma_C$
is a substitution over $C$ with abelianization $N_C$, primitive by
normalization. $\ell$-fixedness says $\Sigma_C(s^{*}_{\ell})$ begins
with $s^{*}_{\ell}$, so the one-sided fixed point $\tau$ exists. Each
cell of the rays' common refinement occurs, suitably translated, in
every sufficiently deep refinement of $s^{*}_{\ell}$ along the tracked
anchoring of Lemma~\ref{lem:tel}, hence lies in $C$, and the level-$n$
cell sequences are by construction $\Sigma_C^{\,n}(s^{*}_{\ell})$;
pass to the limit. The width vector $w$ is the Perron eigenvector of
$N_C$ (Theorem~\ref{thm:floor}), giving the self-similar realization
with inflation $\beta$. The degree statement is
Lemma~\ref{lem:field}; injectivity of $V$ gives the intertwining.
\end{proof}

\begin{remark}[Placement and contact surface]\label{rem:transplace}
Theorem~\ref{thm:transversal} moves the counterexample shape into a
classical-looking category: primitive substitutions with reducible
Pisot--Perron expansion, fibered over an irreducible Pisot
substitution, with a diagonal-avoidance property along the fixed
point. The avoidance itself remains positional --- it is the
cut-disjointness of Theorem~\ref{thm:frozen} re-expressed --- but the
package now sits where the reducible-Pisot and Property~(F) /
bounded-remainder literatures have tools, and where any rigidity
statement for reducible expansions with irreducible Pisot degree would
bear directly on Conjecture~\ref{conj:H3}.
\end{remark}

\begin{lemma}[Unit step]\label{lem:unitstep}
Let $s_1,s_2$ be consecutive cells (left to right) of the common
refinement of two tile grids, with $s_1=(i_1,\delta_1,j_1)$,
$s_2=(i_2,\delta_2,j_2)$, and suppose their common boundary is not a
shared cut of the two grids. Then exactly one letter changes, and
\[
\delta_2-\delta_1=
\begin{cases}
-e_{i_1}, & j_2=j_1 \ \text{(an $I$-grid cut: the $i_1$-tile ends)},\\
+e_{j_1}, & i_2=i_1 \ \text{(a $J$-grid cut: the $j_1$-tile ends)}.
\end{cases}
\]
At a shared cut both letters change and
$\delta_2-\delta_1=e_{j_1}-e_{i_1}$ (the double step), and the
configuration produces aligned cells as in Lemma~\ref{lem:cut}.
\end{lemma}

\begin{proof}
Crossing an $I$-cut increments the prefix-Parikh vector of the $I$-grid
by $e_{i_1}$, which enters the offset with a minus sign; a $J$-cut
enters with a plus sign; both formulas are immediate from
$\delta=A^{\!\top}\delta_{\mathrm{parent}}+\pi(\mathrm{pp}_J)-\pi(\mathrm{pp}_I)$
read across one cut. A simultaneous cut composes the two.
(Machine validation: $1{,}148$ consecutive-cell boundaries across three
certified specimens; the rule holds at every boundary that is not a
shared cut, and all six exceptions decode exactly as shared-cut
configurations adjacent to coincidence cells.)
\end{proof}

\begin{corollary}[The avoidance walk is a substitutive unit-step
walk]\label{cor:walk}
For a trapped component, normalized as in Theorem~\ref{thm:frozen},
the offset sequence $(\delta_n)$ along the transversal fixed point
$\tau$ of Theorem~\ref{thm:transversal} is a walk with steps in
$\{-e_a\}\cup\{+e_b\}$ --- the sign recording which ray's tile ends at
the $n$-th cut, the letter recording which tile --- with no double
steps (cut-disjointness), confined to the closed saturated core
$\mathcal{C}_0$, generated by the primitive substitution $\Sigma_C$,
and avoiding $\{0\}\cup\{e_i-e_j\}$ at every step, forever. The open
core of the conjecture, in its most concrete form to date: such a
substitutive unit-step walk in a finite region with these forbidden
points cannot exist. Its impossibility remains positional --- the walk
simply never lists the forbidden points among its finitely many
visited values --- and is untouched here.
\end{corollary}

\begin{lemma}[Type preservation; fibering over the Barge--Diamond
complex]\label{lem:type}
In a trapped transversal, call a cell boundary \emph{type $I$} or
\emph{type $J$} according to which ray's grid has a vertex there
(exactly one does, by cut-disjointness). The substitution dynamics on
transitions --- sending the boundary between $\tau_n,\tau_{n+1}$ to the
boundary between the last child of $\tau_n$ and the first child of
$\tau_{n+1}$ --- preserves type. On type-$I$ transitions it acts on the
underlying $I$-letter pair by $(i,i')\mapsto(\mathrm{last}\,\sigma(i),
\mathrm{first}\,\sigma(i'))$, i.e.\ it fibers over the classical
Barge--Diamond two-letter dynamics of $\sigma$; symmetrically for
type~$J$.
\end{lemma}

\begin{proof}
The boundary at $x$ maps to the boundary at $\beta x$ under the
self-similarity of the frozen pair. If $x$ is an $I$-vertex, $\beta x$
is a vertex of the refined $I$-grid; it cannot also be a vertex of the
refined $J$-grid, for that would be a shared cut of the pair --- which
is self-similar, hence cut-disjoint at every level
(Theorem~\ref{thm:frozen}). The $I$-letters at the image boundary are
the last letter of $\sigma(i)$ (the refined $I$-tile ending at
$\beta x$) and the first letter of $\sigma(i')$, while the ambient
$J$-letter is unchanged across $\beta x$.
\end{proof}

\begin{remark}[Localization of the homological-Pisot question for
$\Sigma_C$]\label{rem:QA}
By Lemma~\ref{lem:type}, the transition subcomplex of $\Sigma_C$
decomposes edge-wise into an $I$-part and a $J$-part, each carrying
dynamics fibered over the Barge--Diamond complex of $\sigma$ itself
(whose eventual-range structure, for irreducible $\sigma$ without
asymptotic cycles, is the classical loop-free object). Whether
$\Sigma_C$ is homological Pisot --- which by the structural fact of
Barge--Bruin--Jones--Sadun would force every eigenvalue of $N_C$
outside $\chi_A$ to be zero or a root of unity, a spectral shape
realized by no recurrent component in the certified corpus ---
therefore reduces to: (1)~the spectral condition on $N_C$ itself,
(2)~the component count of the eventual range against the multiplicity
of the eigenvalue~$1$, and (3)~loop-freeness of the eventual range,
now a question about cycles in a complex fibered over a known
loop-free base, settled by the orientation/reducedness argument of
Lemma~\ref{lem:redproj}. Section~\ref{sec:trans}
resolves (3) affirmatively over loop-free bases
(Theorem~\ref{thm:inherit}) and converts (2) into an unconditional
spectral constraint (Corollary~\ref{cor:rootofunity}); the residue of
the rigidity question is the spectral condition (1) in the sharpened
form of Corollary~\ref{cor:QAresidue}.
\end{remark}

\begin{remark}\label{rem:powerverif}
Computational sanity (elimination-grade): $\deg\beta^p=3$ verified for
sampled specimens across $p\in\{2,3,6,12,24\}$, and the full PIP
closure of $\sigma^p$ (primitivity, irreducibility, distinctness of
$\beta_m^{\,p}$) for six specimens across $p\in\{2,3,4,6,8,12\}$.
Corollary~\ref{cor:power} explains a feature of the recurrent-cycle
census: the resonance quantity $|1-\beta_2^{\,c}|$ over realized cycle
lengths is bounded below by $0.347$ across $3{,}145$ certified
specimens and never approaches $0$ through angular near-resonance ---
as it cannot exactly, $\arg\beta_2$ being an irrational angle. The
hypothetical counterexample to Conjecture~\ref{conj:H3} is now
simultaneously constrained in size ($|C|\ge p\,k$, and WLOG $p=1$ with
$N_C$ primitive), location (the saturated core of \S\ref{sec:core}),
spectrum ($\rho(N_C)=\beta$ exactly, $\chi_A\mid\chi_{N_C}$), and
geometry (totally cut-disjoint at all scales).
\end{remark}

\section{The transition complex of $\Sigma_C$: decomposition, inherited
loop-freeness, and the coincidence rank of the pair system}\label{sec:trans}

Throughout this section $C$ is a trapped component, normalized as in
Theorem~\ref{thm:transversal}, and --- except in
Lemma~\ref{lem:legfactor} and its corollaries, which need no such
hypothesis --- $C$ is \emph{anchor-free} (the H3$'$-counterexample
shape), so that by Theorem~\ref{thm:frozen} no state has offset $0$ or
$L_i-L_j$ and the transversal is totally cut-disjoint.

We use the exact transition-complex conventions of
Barge--Diamond~\cite{BD08} as stated in
Barge--Bruin--Jones--Sadun~\cite{BBJS12}: for a primitive aperiodic
substitution $\varphi$ on an alphabet $\mathcal B$, the complex $G_0$
has one edge $v_{ab}$ per legal two-letter word $ab$, running from the
vertex $\mathrm{out}_a$ (the end of the letter edge $e_a$) to the
vertex $\mathrm{in}_b$ (the beginning of $e_b$); since the letter
edges are not part of $G_0$, the vertices $\mathrm{out}_a$ and
$\mathrm{in}_a$ are \emph{distinct} and $G_0$ is bipartite.
Substitution acts by $v_{ab}\mapsto
v_{\mathrm{last}\,\varphi(a),\,\mathrm{first}\,\varphi(b)}$ and
permutes the edges of the eventual range $G_0^{ER}$. With $b_1=E-V+C$
the first Betti number of $G_0^{ER}$ and $\mathcal C$ its component
count, the exact sequence of \cite{BD08,BBJS12} (for a power fixing
the $ER$ edges, whose image of $\widetilde H^0(G_0^{ER})$ lies in the
$+1$ eigenspace of the transposed abelianization) gives
$\dim\check H^1(\Omega_\varphi,\mathbb Q)=\dim\varinjlim
M_\varphi^{\!\top}-(\mathcal C-1)+b_1$.

\begin{lemma}[Connectivity for irreducible Pisot]\label{lem:connect}
If $\varphi$ is irreducible Pisot of degree $k$ then $G_0^{ER}(\varphi)$
is connected and $\dim\check H^1(\Omega_\varphi)=k+b_1(G_0^{ER})$;
$\varphi$ is homological Pisot iff $b_1=0$ iff $\varphi$ has no
asymptotic cycles.
\end{lemma}

\begin{proof}
$1$ is not an eigenvalue of any power of the abelianization: the
eigenvalues of the $q$-th power are $\beta^q$ and $\beta_m^{\,q}$ with
$\beta^q>1$ and $|\beta_m|<1$. So $\widetilde H^0(G_0^{ER})$ injects
into the trivial $+1$ eigenspace, forcing $\mathcal C=1$, and
$\dim\varinjlim A^{\!\top}=k$ since $\det A\ne 0$.
\end{proof}

(Machine validation: $\mathcal C=1$ on all $4{,}554$ exhaustive $k=3$
PIP specimens and $1{,}190$ random $k=4$ specimens, exact-convention
instrument gated against the calibration examples of \cite{BBJS12},
including the asymptotic-cycle specimen $1\mapsto 21112$, $2\mapsto
121$ with $b_1=1$ and the coincidence-rank-$3$ triple cover with
$\mathcal C=3$, $b_1=0$.)

\begin{remark}\label{rem:convention}
The integers $\mathcal C$ and $b_1$ are invariants of the pair
$(\Omega,$ chosen Barge--Diamond complex$)$, not of $\Omega$ alone: a
collared presentation can change both while preserving $\check
H^1(\Omega)$ (e.g.\ the uncollared Fibonacci $S_0^{ER}$ is contractible
with $\mathcal C=1$, a collared one has $\mathcal C=2$; both give
$\check H^1=\mathbb Z^2$). We fix the uncollared one-edge-per-legal-
two-word convention of \cite{BD08,BBJS12} throughout, so all reported
$b_1$ are mutually comparable; every statement of this section uses
$\mathcal C$ and $b_1$ of that fixed complex and is set against the
homeomorphism invariant $\dim\check H^1$ (cf.\ the general
$1$-dimensional sequence as recorded in \cite{Sa14}).
\end{remark}

\begin{lemma}[Boundary-type determinacy]\label{lem:typedet}
Let $c=(i,\delta,j)$ be any genuine overlap with $t=\langle\delta,
L\rangle$, so $o(c)=[\max(0,t),\,\min(L_i,t+L_j))$. The left endpoint
of $o(c)$ is a vertex of the $I$-grid iff $t<0$, of the $J$-grid iff
$t>0$, and of both iff $t=0$ (the prefix anchor); the right endpoint
is a vertex of the $I$-grid iff $t>L_i-L_j$, of the $J$-grid iff
$t<L_i-L_j$, and of both iff $t=L_i-L_j$ (the suffix anchor). In
particular, for an anchor-free trapped component every cell has a
well-defined left type and right type, intrinsic to its state, and the
type of the boundary between consecutive transversal cells equals the
right type of the left cell and the left type of the right cell.
\end{lemma}

\begin{proof}
Immediate from $o(c)=[0,L_i)\cap[t,t+L_j)$: the left endpoint is
$0$ (an $I$-tile vertex, interior to the $J$-tile) when $t<0$ and $t$
(a $J$-tile vertex, interior to the $I$-tile) when $t>0$; symmetrically
on the right. By irreducibility, $t=0$ iff $\delta=0$ and $t=L_i-L_j$
iff $\delta=e_i-e_j$, the two anchors. Consistency across a boundary is
Lemma~\ref{lem:unitstep} read geometrically. (Machine validation:
$9{,}945$ consecutive-cell boundaries across $24$ PIP specimens,
exact integer-vector cut positions; the sign predicates match the
geometric type at all $9{,}660$ non-shared boundaries, and all $285$
shared cuts are anchor-adjacent.)
\end{proof}

\begin{theorem}[Type decomposition]\label{thm:decomp}
For an anchor-free trapped $C$, every edge of $G_0(\Sigma_C)$ into
$\mathrm{in}_t$ has type equal to the left type of $t$, and every edge
out of $\mathrm{out}_s$ has type equal to the right type of $s$.
Consequently $G_0(\Sigma_C)$, and hence $G_0^{ER}(\Sigma_C)$, is the
\emph{vertex-disjoint} union of its type-$I$ part $G_I$ and type-$J$
part $G_J$; in particular $b_1(G_0^{ER}(\Sigma_C)) = b_1(G_I) +
b_1(G_J)$. Both parts meet the eventual range, so
\[
\mathcal C\bigl(G_0^{ER}(\Sigma_C)\bigr)\;\ge\;2 .
\]
\end{theorem}

\begin{proof}
The first sentence is Lemma~\ref{lem:typedet}: the type of an edge
$(s,t)$ is the type of the shared boundary, which is determined by
either incident state. A vertex is therefore incident to edges of one
type only, and an edge joins an out-vertex and an in-vertex of its own
type: the two parts share no vertices. Both parts are nonempty in the
eventual range: type-$I$ and type-$J$ legal two-words exist at every
level (each ray's tiles are finite, so each grid cuts the transversal
syndetically), the substitution map preserves type
(Lemma~\ref{lem:type}), and the eventual range is the image of the
full legal set under a high power.
\end{proof}

\begin{corollary}[Root-of-unity spectrum; size-floor upgrade]
\label{cor:rootofunity}
Write $\chi_{N_C}=\chi_A\cdot q(x)$ (Theorem~\ref{thm:floor}). For an
anchor-free trapped $C$, at least $\mathcal C(G_0^{ER}(\Sigma_C))-1\ge
1$ roots of $q$, counted with multiplicity, are roots of unity. In
particular $\deg q\ge1$, so $|C|\ge k+1$, strictly improving
Theorem~\ref{thm:floor}; and $\operatorname{spec}(N_C)$ must contain a
root of unity. Containment of \emph{a} root of unity is not by itself
rare --- it occurs in $225$ of $358$ productive recurrent components of
the certified corpus. The trapped fingerprint is the \emph{exhaustive}
clause of Corollary~\ref{cor:QAresidue}, that \emph{every} nonzero root
of $q$ be a root of unity (equivalently $\dim\check
H^1(\Omega_{\Sigma_C})=k$); this is realized by no productive
component, $226$ of $245$ sampled carrying an off-unit-circle Perron
root $\rho(N_C)>1$.
\end{corollary}

\begin{proof}
In the exact sequence for the edge-fixing power $\Sigma_C^{\,q}$, the
group $\widetilde H^0(G_0^{ER})\cong\mathbb Q^{\mathcal C-1}$ injects
into the fixed space of the limit automorphism, so the eigenvalue $1$
of $N_C^{\,q}$ has geometric multiplicity at least $\mathcal C-1$;
each such eigenvalue comes from a root $\zeta$ of $\chi_{N_C}$ with
$\zeta^q=1$, and no root of $\chi_A$ is a root of unity (Pisot).
Theorem~\ref{thm:decomp} gives $\mathcal C\ge2$.
\end{proof}

\begin{lemma}[Reduced projection]\label{lem:redproj}
Every simple cycle of $G_0^{ER}(\Sigma_C)$ is type-pure
(Theorem~\ref{thm:decomp}) and projects, via the letter coordinate of
its type ($s\mapsto i(s)$ for type $I$, $s\mapsto j(s)$ for type $J$),
to a \emph{reduced} closed edge-path in $G_0^{ER}(\sigma)$; hence its
existence forces $b_1(G_0^{ER}(\sigma))\ge1$.
\end{lemma}

\begin{proof}
Take the cycle type-$I$ (type $J$ is symmetric in the other
coordinate). A type-$I$ edge $(s,t)$ projects to the legal
$\sigma$-two-word $(i(s),i(t))$ of consecutive $I$-tiles; by
Lemma~\ref{lem:type} the projection intertwines the two-word
dynamics, and since $ER$ edges lie on cycles of the edge permutation,
the projected edge lies on a cycle of $\sigma$'s two-word map, hence
in $G_0^{ER}(\sigma)$. By Lemma~\ref{lem:unitstep}, a type-$I$
predecessor of $t$ is determined by its $i$-letter
($s=(i(s),\,\delta_t+e_{i(s)},\,j(t))$) and a type-$I$ successor of
$s$ by its $i$-letter; so two distinct cycle edges sharing
$\mathrm{in}_t$ (resp.\ $\mathrm{out}_s$) project to \emph{distinct}
base edges sharing $\mathrm{in}_{i(t)}$ (resp.\
$\mathrm{out}_{i(s)}$). Consecutive edges of the alternating chain
always share an $\mathrm{in}$- or $\mathrm{out}$-vertex, so the
projected closed walk never traverses a base edge and immediately
retraverses it: it is reduced in the free-group sense (the only
relevant one, since edge reduction is a local, consecutive-pair
operation; non-consecutive coincidences of projected vertices are
irrelevant). Each type-$I$ edge advances the $i$-coordinate by one
$I$-tile (Lemma~\ref{lem:unitstep}), so the walk has positive length.
A reduced nontrivial closed edge-walk represents a nontrivial element
of $\pi_1(G_0^{ER}(\sigma))$; a graph with nontrivial fundamental
group is not a tree, so $b_1(G_0^{ER}(\sigma))\ge1$. (The argument is
through $\pi_1$, not the homology class of the projected loop itself,
which may vanish in $H_1$; existence of \emph{any} reduced nontrivial
closed walk already forces a cycle in the graph.)
\end{proof}

\begin{theorem}[Inherited loop-freeness]\label{thm:inherit}
If $\sigma$ has no asymptotic cycles ($b_1(G_0^{ER}(\sigma))=0$;
equivalently, by Lemma~\ref{lem:connect}, $\dim\check
H^1(\Omega_\sigma)=k$), then for every anchor-free trapped component
$C$ of $\sigma$,
\[
b_1\bigl(G_0^{ER}(\Sigma_C)\bigr)=0 .
\]
Condition (3) of Remark~\ref{rem:QA} is therefore automatic over
loop-free bases.
\end{theorem}

\begin{proof}
$b_1>0$ produces a simple cycle; Lemma~\ref{lem:redproj} contradicts
$b_1(G_0^{ER}(\sigma))=0$.
\end{proof}

\begin{corollary}[The Q-A residue]\label{cor:QAresidue}
Over a loop-free base, $\Sigma_C$ is homological Pisot if and only if
the number of nonzero roots of $q$, with multiplicity, equals
$\mathcal C(G_0^{ER}(\Sigma_C))-1$; in that case every nonzero root of
$q$ is a root of unity. The rigidity question for the trapped shape is
thus exactly a counting identity between the extra spectrum of $N_C$
and the components of its transition complex.
\end{corollary}

\begin{remark}[The residue is irreducible to the base loop-freeness]
\label{rem:cond1-irreducible}
On the loop-free base stratum Theorem~\ref{thm:inherit} supplies
$b_1(G_0^{ER}(\Sigma_C))=0$, and Corollary~\ref{cor:rootofunity}
supplies the lower half of Corollary~\ref{cor:QAresidue}: at least
$\mathcal C-1$ of the nonzero roots of $q$ are roots of unity. The
entire residue is then the matching \emph{upper} bound, that the number
$r$ of nonzero roots of $q$ does not exceed $\mathcal C-1$. This upper
bound is not a consequence of loop-freeness: for a producing within-SCC
matrix the nonzero spectrum is essentially full (across $152$ recurrent
components the ratio of nonzero spectrum to size has median $1.000$,
and only $7.9\%$ have at most $k-1$ nonzero eigenvalues), so
$r\le\mathcal C-1$ is a strongly non-generic spectral collapse ---
precisely the homological-Pisot hypothesis itself, not a corollary of
$b_1=0$. Equivalently, for connected $G_0^{ER}(\Sigma_C)$ it demands
$q(x)=x^{|C|-k}$ (all extra spectrum nilpotent), the spectrally minimal
trapped shape. The identity $r=(\mathcal C-1)-b_1$ and the
root-of-unity clause are confirmed on the two homological-Pisot triple
covers of \cite{BBJS12}: there $q=x^{a}(x-1)^2$ with $\mathcal C=3$,
$b_1=0$, $r=2$.
\end{remark}

\begin{proof}
$\dim\check H^1(\Omega_{\Sigma_C})
=\dim\varinjlim N_C^{\!\top}-(\mathcal C-1)+b_1
=\bigl(k+\#\{\text{nonzero roots of }q\}\bigr)-(\mathcal C-1)$ by
Theorem~\ref{thm:inherit}; set this equal to $k=\deg\beta$. Equality
forces all nonzero roots of $q$ to be roots of unity by
Corollary~\ref{cor:rootofunity}.
\end{proof}

\begin{lemma}[Leg factor]\label{lem:legfactor}
For any trapped component $C$ (no anchor-freeness needed),
$(\Omega_{\Sigma_C},\mathbb R)$ factors onto $(\Omega_\sigma,\mathbb
R)$.
\end{lemma}

\begin{proof}
A tiling $S\in\Omega_{\Sigma_C}$ assigns to $\mathbb R$ a bi-infinite
cell sequence with widths $w(s)$; merge maximal runs of cells between
consecutive $I$-grid vertices (well defined by
Lemma~\ref{lem:typedet}, shared or not) and label each merged tile by
the common $i$-letter. $I$-vertices are syndetic (each $I$-tile has
length at most $L_{\max}$ and cells have width at least $\min_s
w(s)>0$), every patch of $S$ is a translate of a patch of $\tau$,
where the merged tiles are precisely the $I$-grid tiles of the frozen
ray of length $L_{i}$, so the merged object is locally allowed for
$\sigma$ and $\pi_1(S)\in\Omega_\sigma$. The local rule $\pi_1$ is
continuous and translation-equivariant; its image is a nonempty closed
invariant subset of the minimal $\Omega_\sigma$, hence everything.
\end{proof}

\begin{corollary}[Coincidence rank of the pair system]\label{cor:cr}
$\mathrm{cr}(\Sigma_C)\ge2$ for every trapped component. If moreover
$\sigma$ is unimodular and $\Sigma_C$ is homological Pisot, then
$\mathrm{cr}(\Sigma_C)\ge3$; and if $\mathrm{cr}(\Sigma_C)=2$ (for
unimodular $\sigma$) then $\dim\check H^1(\Omega_{\Sigma_C})\ge2k-1$.
\end{corollary}

\begin{proof}
$\Sigma_C$ is a primitive, non-periodic, FLC substitution on letters
with Pisot dilatation $\beta$ (flow-periodicity of a transversal would
push to flow-periodicity of $\pi_1(S)$, impossible), so
$\mathrm{cr}(\Sigma_C)=1$ iff $\Omega_{\Sigma_C}$ has pure discrete
spectrum~\cite[Thm.~4(5)]{B15}. Pure discreteness passes to factors,
and $\Omega_\sigma$ is a factor (Lemma~\ref{lem:legfactor}); but
$\sigma$ carries a trapped component, hence fails overlap coincidence
and pure discreteness~\cite{Sol97,AL11}. So $\mathrm{cr}(\Sigma_C)\ge
2$. (The trapped component here is the \emph{overlap-graph} object, for which
overlap coincidence $\Leftrightarrow$ pure discreteness is a genuine equivalence.
Its identification with the balanced-pair closed nonproductive SCC of the main
manuscript is the ``bridge converse'' made explicit there (PSC\_PROOF v15,
\S6.5): balanced-pair non-termination implies overlap non-coincidence only through
that identification, which the main manuscript therefore states as a hypothesis.) The norm of $\beta$ is $\pm\det A=\pm1$, odd; if
$\mathrm{cr}(\Sigma_C)=2$ then $\dim\check
H^1(\Omega_{\Sigma_C})\ge2k-1$ by~\cite[Thm.~22]{B15} (whose
hypotheses are exactly: Pisot substitution on letters, odd norm,
coincidence rank two --- reducibility of $\Sigma_C$ is permitted),
which exceeds $k$ for $k\ge2$, contradicting homological-Pisotness.
\end{proof}

\begin{theorem}[Conditional placement: the unimodular conjecture in the
topological orbit]\label{thm:QACRC}
Let Q-A denote the statement that $\Sigma_C$ is homological Pisot for
every trapped component of every unimodular irreducible Pisot
substitution, and CRC the Coincidence Rank Conjecture
of~\cite{BBJS12,B15}. Then
\[
\text{Q-A}\;+\;\text{CRC}\;\Longrightarrow\;\text{unimodular PSC}.
\]
\end{theorem}

\begin{proof}
Suppose $\sigma$ is unimodular irreducible Pisot without pure discrete
spectrum. Failure of overlap coincidence produces a closed recurrent
noncoincident component with $\rho(N_C)=\beta$ and $\mathrm{leak}
\equiv0$ (Theorem~\ref{thm:dichotomy} with the Akiyama--Lee criterion
\cite{AL11}; cf.\ Remark~\ref{rem:calib}), i.e.\ a trapped component,
which persists under the power normalization of
Theorem~\ref{thm:powerred} (pure discreteness is power-invariant). By
Q-A, $\dim\check H^1(\Omega_{\Sigma_C})=\deg\beta$, and the norm of
$\beta$ is $\pm1$, so CRC gives $\mathrm{cr}(\Sigma_C)\mid 1$, i.e.\
pure discreteness of $\Omega_{\Sigma_C}$~\cite[Thm.~4(5)]{B15}, hence
of its factor $\Omega_\sigma$ (Lemma~\ref{lem:legfactor}) ---
contradicting the hypothesis.
\end{proof}

\begin{remark}[Honest placement; adjudication of the branching
question]\label{rem:adjud}
(i) CRC is open precisely in the regime $\mathrm{cr}\ge3$ that
Corollary~\ref{cor:cr} lands in (it is proved for degree one
\cite{BBJS12} and for coincidence rank two \cite{B15}), and Q-A is
open; Theorem~\ref{thm:QACRC} transfers the unimodular conjecture into
the orbit of the homological-Pisot/CRC program rather than resolving
anything. (ii) All results of this section are count- and
topology-level constraints on the hypothetical trapped shape; the
positional avoidance residue of Corollary~\ref{cor:walk} is untouched,
as the no-coarse-quotient obstruction requires. (iii) A formulation of
the loop question as ``backward branching of the avoidance walk''
(two trapped pasts of one tail) has no content: asymptotic pairs exist
in every primitive aperiodic substitution subshift, every point of the
hull of $\tau$ decodes to a trapped pair, and recognizability of
$\Sigma_C$ constrains desubstitution, not backward determinism of the
shift; the refutable object is exactly $b_1(G_0^{ER}(\Sigma_C))$ ---
by Theorem~\ref{thm:decomp} a type-pure alternating cycle of
asymptotic branchings --- and over loop-free bases it is now excluded.
(iv) Barge's dichotomy \cite[Cor.~25]{B15} (a PSC counterexample has
an asymptotic cycle or coincidence rank at least three) calibrates the
base-level census: in the exhaustive $k=3$, $|\sigma(a)|\le3$ corpus,
$3{,}102$ of $4{,}554$ specimens (including $1{,}740$ of the $2{,}628$
unimodular ones) are loop-free, hence homological Pisot with the
$\mathrm{cr}\ge3$ alternative; the $b_1>0$ stratum, exactly
quantified ($b_1\in\{1,2,4\}$ at $k=3$; $b_1$ up to the complete
bipartite ceiling at $k=4$, where $b_1=3$ also occurs), is where the
asymptotic-cycle alternative lives. As always, no census bears on the
uniform statements.
\end{remark}

\section{Computational evidence (elimination, not validation)}

All computations use the validated instrument
(\texttt{overlap\_residual.py}; partition identity verified to
$5.6\times10^{-16}$, per-SCC restriction to $4.8\times10^{-14}$, gap
formula to $2.1\times10^{-14}$ including an SCC of size $5068$).

\begin{center}
\begin{tabular}{lrrrr}
\toprule
corpus & closed SCCs & nonproducers & H3 violators & notes\\
\midrule
$k=3$, $|\sigma(a)|\le3$, exhaustive (4{,}554 PIP) & 1{,}386 & 0 & 0 &
zero-free $\Rightarrow$ suffix-aligned, $324/324$\\
$k=3$, $|\sigma(a)|\le4$, random & 105 & 0 & 0 & \\
$k=3$, $|\sigma(a)|\le6$, random & 35 & 0 & 0 & 12 capped excluded\\
$k=4$, $|\sigma(a)|\le4$, random & $\sim$60 & 0 & 0 & \\
$k=5$, $|\sigma(a)|\le3$, random & 12 & 0 & 0 & \\
\bottomrule
\end{tabular}
\end{center}

$Q(p)$ census (abstract states; no reachability or pure-discreteness
input): collisions in $3{,}216/3{,}216$ cases over $p\in\{1,2,3,6\}$,
with maximal cone depth $6,3,3,1$ respectively --- the small-offset
regime $t^*=c/(\beta^p-1)\to0^+$ is the easy one, and the difficulty of
$Q(p)$ concentrates at small $p$ (sampled pattern, $k=3$).

Both anchor-only classes occur (prefix-only and suffix-only closed SCCs
exist), so the disjunction in Theorem~\ref{thm:C} cannot be strengthened
to a single anchor type. By Theorem~\ref{thm:C} the anchoring counts in
the table are implied by the production counts; they are therefore
consistency checks, and \emph{no census evidence bears directly on
Conjecture~\ref{conj:H3}} --- necessarily, since a trapped specimen
would refute the Pisot conjecture. Sampled minima/maxima are not proved
bounds; capped specimens are inconclusive and excluded; image length is
a cap throughout.

\begin{thebibliography}{99}
\bibitem{AI01} P.~Arnoux, S.~Ito, \emph{Pisot substitutions and Rauzy
fractals}, Bull. Belg. Math. Soc. Simon Stevin 8 (2001), 181--207.
\bibitem{Ak16} S.~Akiyama, \emph{Strong coincidence and overlap
coincidence}, Discrete Contin. Dyn. Syst. 36 (2016), 5223--5230.
\bibitem{AL11} S.~Akiyama, J.-Y.~Lee, \emph{Algorithm for determining
pure pointedness of self-affine tilings}, Adv. Math. 226 (2011),
2855--2883.
\bibitem{AL14} S.~Akiyama, J.-Y.~Lee, \emph{Overlap coincidence to
strong coincidence in substitution tiling dynamics}, European J.
Combin. 39 (2014), 233--243.
\bibitem{DT89} J.-M.~Dumont, A.~Thomas, \emph{Syst\`emes de
num\'eration et fonctions fractales relatifs aux substitutions},
Theoret. Comput. Sci. 65 (1989), 153--169.
\bibitem{Lee07} J.-Y.~Lee, \emph{Substitution Delone sets with pure
point spectrum are inter-model sets}, J. Geom. Phys. 57 (2007),
2263--2285.
\bibitem{BD02} M.~Barge, B.~Diamond, \emph{Coincidence for substitutions
of Pisot type}, Bull. Soc. Math. France 130 (2002), 619--626.
\bibitem{BK06} M.~Barge, J.~Kwapisz, \emph{Geometric theory of
unimodular Pisot substitutions}, Amer. J. Math. 128 (2006), 1219--1282.
\bibitem{Sing06} B.~Sing, \emph{Pisot substitutions and beyond}, PhD
thesis, Univ. Bielefeld, 2006.
\bibitem{Sol97} B.~Solomyak, \emph{Dynamics of self-similar tilings},
Ergodic Theory Dynam. Systems 17 (1997), 695--738.
\bibitem{BD08} M.~Barge, B.~Diamond, \emph{Cohomology in one-dimensional
substitution tiling spaces}, Proc. Amer. Math. Soc. 136 (2008),
2183--2191.
\bibitem{BBJS12} M.~Barge, H.~Bruin, L.~Jones, L.~Sadun,
\emph{Homological Pisot substitutions and exact regularity}, Israel J.
Math. 188 (2012), 281--300.
\bibitem{B15} M.~Barge, \emph{Factors of Pisot tiling spaces and the
Coincidence Rank Conjecture}, Bull. Soc. Math. France 143 (2015),
357--381; arXiv:1301.7094.
\bibitem{Sa14} L.~Sadun, \emph{Cohomology of hierarchical tilings},
in \emph{Mathematics of Aperiodic Order}, Progr. Math. 309,
Birkh\"auser (2015), 73--104; arXiv:1406.0882.
\end{thebibliography}
\end{document}
